% ============================================
% Section: Implementation
% ============================================
\label{sec:implementation}

Co-LIC2 is implemented as an extension of the open-source Gaussian-LIC2
codebase~\cite{lang2026gaussian2}, preserving full backward compatibility.

\paragraph{Build and dependencies.} The per-node front-end reuses
Gaussian-LIC2's dependency stack: ROS~2 Humble, OpenCV 4.7 (CUDA), LibTorch,
TensorRT 8.6, CUDA 11.7, and the Coco-LIC B-spline library. Cooperative
components add: gRPC (cross-subnet fallback), tinycbor (CBOR encoding), zstd
(entropy coding), GTSAM (pose-graph optimization), DBoW2/NetVLAD (visual loop
closure), and RaptorQ for keyframe FEC. All additions are gated behind CMake
option \texttt{COOP\_ENABLED=ON}; a vanilla build remains faithful
Gaussian-LIC2.

\paragraph{Component layout.} New components reside under \texttt{src/coop/}:
\begin{itemize}
  \item \texttt{coordinator} --- node discovery, descriptor exchange, cluster-head
    election (Alg.~\ref{alg:election});
  \item \texttt{network\_monitor} --- per-link throughput/PDR estimation
    (Eqs.~\ref{eq:channel}) and routing table maintenance;
  \item \texttt{qos\_classifier} --- message-to-stream mapping with EDCA
    DSCP marking;
  \item \texttt{scheduler} --- deadline-aware submap pacing
    (Alg.~\ref{alg:scheduler});
  \item \texttt{submap\_extractor} --- spatial sub-volume carving from
    octree-accelerated Gaussian map;
  \item \texttt{submap\_compressor} --- knob grid search (Eq.~\ref{eq:opt}) +
    budget controller feedback loop;
  \item \texttt{comm\_layer} --- DDS/gRPC transport with QoS profiles;
  \item \texttt{map\_merger} --- ICP registration (Sec.~\ref{sec:registration})
    + Gaussian fusion (Sec.~\ref{sec:merge}) + conflict resolution
    (Eq.~\ref{eq:confidence});
  \item \texttt{loop\_closure\_agent} --- geometric + visual loop detection;
  \item \texttt{pose\_graph} --- GTSAM iSAM2 wrapper with DCS robust kernel.
\end{itemize}
Each component is a ROS~2 LifecycleNode with documented interface contracts.
The generalized multi-sensor front-end modifies only Coco-LIC's
\texttt{spline\_estimator} to parameterize $N_c$, $N_l$.

\paragraph{Configuration.} All parameters (Table~\ref{tab:config}) are ROS~2
parameters, tunable per deployment via YAML without recompilation.

\begin{table}[t]
  \centering
  \caption{Key configuration parameters.}
  \label{tab:config}
  \footnotesize
  \setlength{\tabcolsep}{3pt}
  \begin{tabular}{@{}p{2.2cm}p{1.1cm}p{4.5cm}@{}}
    \toprule
    Parameter & Default & Description \\
    \midrule
    $S_{\text{sub}}$       & 5 m    & submap spatial period \\
    $T_{\text{sub}}$       & 3 s    & submap temporal period \\
    $\Bbudget$             & 20 Mbps & per-node uplink budget \\
    $\tau_d$               & $3\bar{s}$ & Gaussian dedup radius \\
    $T_{\text{heartbeat}}$ & 200 ms & heartbeat interval \\
    $T_{\text{head\_to}}$  & 3 s    & head-loss timeout \\
    $T_{\text{keep}}$      & 60 s   & orphan submap GC delay \\
    $T_{\text{lease}}$     & 1 s    & minimum head tenure \\
    $d_{\max}$             & 2      & max SH degree on wire \\
    $\Delta q$             & 2 cm   & default spatial quantize \\
    $l_{\text{zstd}}$      & 5      & default zstd level \\
    $\Npv_{\min}$          & 5000   & minimum NPV threshold \\
    \bottomrule
  \end{tabular}
\end{table}

\paragraph{Backward compatibility.} Setting \texttt{coop\_enabled:=false} runs
the unmodified Gaussian-LIC2 pipeline---identical odometry, identical local
map PSNR---satisfying the zero-regression goal on all existing single-node
datasets (FAST-LIVO2, M2DGR, MCD).

\paragraph{Testing plan.} We plan a three-phase validation: (1) unit-test
suite for protocol FSMs, serialization, and scheduler convergence (in simulation);
(2) multi-rig emulation on a single machine with virtualized wireless links
(ns-3 integration for Wi-Fi~6 channel models); (3) physical deployment with
2--4 heterogeneous rigs (LiDAR+stereo UGV + monocular-inertial drone + surveying
vehicle) over real Wi-Fi~6 APs, measuring all metrics defined in
Section~\ref{sec:analysis}.
